\documentclass[11pt]{article}
\usepackage[margin=1in]{geometry}
\usepackage{enumitem}
% =============================================================================
% Preambles/header.tex — shared LaTeX/Beamer preamble for this project
%
% Use from a Beamer lecture by prepending:
%     \input{header}
% and compiling with TEXINPUTS=../Preambles:$TEXINPUTS (the /compile-latex
% skill does this automatically).
%
% PALETTE CONTRACT. Color names defined here MUST match the names in
% Quarto/theme-template.scss so Beamer and Quarto renderings use the same
% palette. The scripts/check-palette-sync.sh script enforces this.
%
% To customize: edit the HEX values below AND the matching $variable in
% Quarto/theme-template.scss, then run scripts/check-palette-sync.sh.
% =============================================================================

% -----------------------------------------------------------------------------
% Engine detection + encoding
%
% Our /compile-latex skill uses XeLaTeX, where inputenc is unnecessary and
% can produce encoding warnings. Only load inputenc under pdfTeX.
% -----------------------------------------------------------------------------
\usepackage{iftex}
\ifPDFTeX
  \usepackage[utf8]{inputenc}
\fi

\usepackage{amsmath, amssymb, amsthm}
\usepackage{booktabs}
\usepackage{graphicx}

% -----------------------------------------------------------------------------
% PALETTE — mirrors Quarto/theme-template.scss variable names.
% Load xcolor and define colors BEFORE anything that references them
% (notably hyperref, hypersetup, and the Beamer color assignments).
% -----------------------------------------------------------------------------
\usepackage{xcolor}

\definecolor{primary-blue}{HTML}{012169}      % headings, accents
\definecolor{primary-gold}{HTML}{B9975B}      % emphasis, borders
\definecolor{highlight-yellow}{HTML}{F2A900}  % markers, alerts
\definecolor{light-bg}{HTML}{E8EDF5}          % title / section backgrounds
\definecolor{jet}{HTML}{1A1A1A}               % body text

% Semantic colors (used in diagrams and annotations)
\definecolor{positive}{HTML}{15803D}          % good / observed / identified
\definecolor{negative}{HTML}{B91C1C}          % bad / problematic / bias
\definecolor{neutral}{HTML}{525252}           % reference / context

% Highlight accents (match .hi-* classes in the SCSS)
\definecolor{hi-slate}{HTML}{314F4F}
\definecolor{hi-green}{HTML}{2E8B57}
\definecolor{hi-red}{HTML}{C41E3A}

% Semantic aliases so skill templates and snippets can write
% \textcolor{accent}{...} without hard-coding "primary-blue".
\colorlet{accent}{primary-blue}
\colorlet{accent2}{primary-gold}

% -----------------------------------------------------------------------------
% Hyperref — loaded AFTER xcolor + palette so link colors resolve.
% -----------------------------------------------------------------------------
\usepackage{hyperref}
\hypersetup{colorlinks=true, linkcolor=primary-blue, urlcolor=primary-blue,
            citecolor=primary-blue, pdfborder={0 0 0}}

% -----------------------------------------------------------------------------
% Course code — set once per deck (after \input{header}, before \begin{document})
% via \coursecode{CS401}. Shown in the footer on every slide so a deck's course
% is identifiable without opening the title page. Empty by default.
% -----------------------------------------------------------------------------
\makeatletter
\newcommand{\coursecode}[1]{\gdef\@coursecodetext{#1}}
\gdef\@coursecodetext{}
\makeatother

% -----------------------------------------------------------------------------
% Beamer theme assignments (only apply under Beamer).
%
% We detect Beamer via \@ifclassloaded{beamer}, which is the documented,
% non-internal way. This must be wrapped in \makeatletter/\makeatother
% because \@ifclassloaded contains an @-catcode character.
% -----------------------------------------------------------------------------
\makeatletter
\@ifclassloaded{beamer}{%
  \usefonttheme{professionalfonts}
  \setbeamercolor{structure}{fg=primary-blue}
  \setbeamercolor{frametitle}{fg=primary-blue}
  \setbeamercolor{title}{fg=primary-blue}
  \setbeamercolor{subtitle}{fg=primary-gold}
  \setbeamercolor{author}{fg=primary-blue}
  \setbeamercolor{institute}{fg=neutral}
  \setbeamercolor{date}{fg=primary-gold}
  \setbeamercolor{itemize item}{fg=primary-blue}
  \setbeamercolor{itemize subitem}{fg=primary-gold}
  \setbeamercolor{alerted text}{fg=primary-gold}
  \setbeamercolor{block title}{fg=white, bg=primary-blue}
  \setbeamercolor{block body}{bg=light-bg}
  % Semantic block variants: block=definition/concept (blue), exampleblock=
  % worked example (green/positive), alertblock=key takeaway or caution (gold).
  % Keep to <=2 colored boxes per slide across ALL three variants (INV-7).
  \setbeamercolor{block title example}{fg=white, bg=positive}
  \setbeamercolor{block body example}{bg=light-bg}
  \setbeamercolor{block title alerted}{fg=white, bg=primary-gold}
  \setbeamercolor{block body alerted}{bg=light-bg}

  % Minimal footer: course code (left) + page X / N (right), no navigation clutter
  \setbeamertemplate{navigation symbols}{}
  \setbeamertemplate{footline}{%
    \color{neutral}\footnotesize\hspace{0.7em}\@coursecodetext\hfill
    \insertframenumber/\inserttotalframenumber\hspace{0.7em}\vspace{0.3em}}
}{}
\makeatother

% -----------------------------------------------------------------------------
% TikZ setup — libraries the snippet gallery uses
% -----------------------------------------------------------------------------
\usepackage{tikz}
\usetikzlibrary{arrows.meta, positioning, calc, decorations.pathreplacing,
                fit, shapes.geometric, backgrounds}

% Shared TikZ styles — snippets and hand-written diagrams can pick these up
% via `\begin{tikzpicture}[every node/.style={...}]` or by naming specific
% styles (e.g., `\node[dag-node] (X) at (0,0) {X};`).
\tikzset{
  % Node styles for causal DAGs and similar
  dag-node/.style = {
    draw=primary-blue, fill=light-bg, rounded corners=2pt,
    minimum width=1.2cm, minimum height=0.9cm, align=center, font=\small
  },
  decision-node/.style = {
    draw=primary-gold, fill=white, diamond, aspect=1.8,
    minimum width=1.8cm, minimum height=1.2cm, align=center, font=\small,
    inner sep=1pt
  },
  % Edge styles — observed vs counterfactual
  observed-edge/.style = {-{Latex[length=2.5mm]}, thick, draw=jet},
  counterfactual-edge/.style = {-{Latex[length=2.5mm]}, thick, draw=neutral, dashed},
  confound-edge/.style = {-{Latex[length=2.5mm]}, thick, draw=negative, dashed},
  % Data-point styles for plots
  observed-dot/.style = {fill=primary-blue, draw=primary-blue, circle, inner sep=1.2pt},
  counterfactual-dot/.style = {fill=white, draw=neutral, circle, inner sep=1.2pt}
}

% -----------------------------------------------------------------------------
% Convenience macros
% -----------------------------------------------------------------------------
\newcommand{\muted}[1]{\textcolor{neutral}{#1}}
\newcommand{\key}[1]{\textcolor{primary-gold}{\textbf{#1}}}
\newcommand{\good}[1]{\textcolor{positive}{#1}}
\newcommand{\bad}[1]{\textcolor{negative}{#1}}

% Standout frame macro: full-bleed dark background for section transitions.
% Beamer-only (uses \begin{frame}/\setbeamercolor/\usebeamerfont) -- do not
% call from an article-class file (e.g. Notes/<CODE>/*-notes.tex). \muted,
% \key, \good, \bad, and \coursecode above are plain xcolor/gdef and are
% safe to use from either class; verified when Notes/article-class support
% was added.
% Expand: \transitionslide{Your transition title}
\newcommand{\transitionslide}[1]{%
  \begingroup
  \setbeamercolor{background canvas}{bg=primary-blue}
  \begin{frame}[plain]
    \centering
    \vfill
    {\usebeamerfont{title}\color{white}\Large #1\par}
    \vfill
  \end{frame}
  \endgroup
}

% Section divider frame (Beamer-only), an alternative to \transitionslide with a
% darker two-line style: near-black (jet) background, a small white label line
% above a larger gold title, and a thin gold rule along the bottom edge.
% Expand: \sectiondivider{Section 2}{Pointers}
\newcommand{\sectiondivider}[2]{%
  \begingroup
  \setbeamercolor{background canvas}{bg=jet}
  \begin{frame}[plain]
    \vspace*{3.0cm}
    {\color{white}\ttfamily\large #1\par}
    \vspace{0.5cm}
    {\color{primary-gold}\LARGE #2\par}
    \begin{tikzpicture}[remember picture, overlay]
      \draw[primary-gold, line width=2.5pt]
        ($(current page.south west)+(0,0.1)$) -- ($(current page.south east)+(0,0.1)$);
    \end{tikzpicture}
  \end{frame}
  \endgroup
}

% -----------------------------------------------------------------------------
% Channel watermark — "Concepts That Click" (YouTube).
%
% Article-class files (CompetitiveExam Q/A sets, Notes, Assignments) get a
% light diagonal draft watermark on every page plus a centered footer naming
% the channel. Beamer decks get a subtle TikZ background watermark instead
% (the footline already carries the course code). Centralized here so every
% MCQ PDF is branded without per-file edits.
% -----------------------------------------------------------------------------
\makeatletter
\@ifclassloaded{article}{%
  \usepackage{draftwatermark}
  \SetWatermarkText{Concepts That Click}
  \SetWatermarkScale{1}
  \SetWatermarkFontSize{2cm}
  \SetWatermarkLightness{0.86}
  \SetWatermarkAngle{45}
  \usepackage{fancyhdr}
  \pagestyle{fancy}
  \fancyhf{}
  \fancyfoot[C]{\footnotesize\textcolor{neutral}{Concepts That Click~|~YouTube: \href{https://www.youtube.com/@concepts.thatclick}{@concepts.thatclick}}}
  \renewcommand{\headrulewidth}{0pt}
  \renewcommand{\footrulewidth}{0pt}
  \fancypagestyle{plain}{%
    \fancyhf{}%
    \fancyfoot[C]{\footnotesize\textcolor{neutral}{Concepts That Click~|~YouTube: \href{https://www.youtube.com/@concepts.thatclick}{@concepts.thatclick}}}%
    \renewcommand{\headrulewidth}{0pt}%
    \renewcommand{\footrulewidth}{0pt}%
  }
}{}
\@ifclassloaded{beamer}{%
  \setbeamertemplate{background}{%
    \begin{tikzpicture}[remember picture, overlay]
      \node[rotate=45, opacity=0.055, align=center]
        at (current page.center)
        {\Huge\textcolor{neutral}{Concepts That Click}};
    \end{tikzpicture}%
  }
}{}
\makeatother 
\coursecode{JKSSB}
\title{Lesson 58 Practice: PYQ-Calibrated Final Retrieval and Audit}
\author{JKSSB Computer Awareness}
\date{\today}
\begin{document}
\maketitle
\noindent\textbf{Provenance:} All 20 questions are original JKSSB-pattern
questions (\textit{[Original]}); none reproduces an official paper item.
Depth mix: L1 recall (Q1--Q10), L2 contrast/classification (Q11--Q17),
L3 one-step application (Q18--Q20); formats include NOT/EXCEPT, matching,
abbreviation, statement-evaluation, and assertion--reason.

\begin{enumerate}[label=\textbf{Q\arabic*.}, leftmargin=*]
\item \textit{[Original]} Expand ISP.
  \begin{enumerate}[label=\Alph*.]
    \item Internet Service Provider
    \item Intranet System Program
    \item Internet Server Process
    \item Internal Service Protocol
  \end{enumerate}
\item \textit{[Original]} Which is the smallest unit of data?
  \begin{enumerate}[label=\Alph*.]
    \item Byte
    \item Nibble
    \item Bit
    \item Word
  \end{enumerate}
\item \textit{[Original]} Pressing which key starts a PowerPoint slide show
  from the beginning?
  \begin{enumerate}[label=\Alph*.]
    \item F7
    \item F9
    \item F5
    \item F1
  \end{enumerate}
\item \textit{[Original]} Every Excel formula must start with which symbol?
  \begin{enumerate}[label=\Alph*.]
    \item +
    \item =
    \item \%
    \item \#
  \end{enumerate}
\item \textit{[Original]} Which technology reads characters printed in
  magnetic ink?
  \begin{enumerate}[label=\Alph*.]
    \item OCR
    \item OMR
    \item MICR
    \item QR code
  \end{enumerate}
\item \textit{[Original]} Integrated circuits (ICs) were introduced in which
  computer generation?
  \begin{enumerate}[label=\Alph*.]
    \item First
    \item Second
    \item Third
    \item Fifth
  \end{enumerate}
\item \textit{[Original]} Which is an open-source mobile operating system?
  \begin{enumerate}[label=\Alph*.]
    \item Android
    \item A proprietary desktop system
    \item A shareware trial system
    \item A firmware loader
  \end{enumerate}
\item \textit{[Original]} Which firmware initialises the hardware and loads
  the operating system at startup?
  \begin{enumerate}[label=\Alph*.]
    \item Device driver
    \item BIOS
    \item Compiler
    \item Spreadsheet
  \end{enumerate}
\item \textit{[Original]} In an email address, the symbol @ separates which
  two parts?
  \begin{enumerate}[label=\Alph*.]
    \item Domain name and host name
    \item User name and domain name
    \item Sender name and recipient name
    \item Subject and body
  \end{enumerate}
\item \textit{[Original]} Which Access query retrieves data on a condition
  without changing stored records?
  \begin{enumerate}[label=\Alph*.]
    \item Update query
    \item Append query
    \item Delete query
    \item Select query
  \end{enumerate}
\item \textit{[Original]} Which of the following is \textbf{NOT} true of
  cache memory?
  \begin{enumerate}[label=\Alph*.]
    \item It holds recently and frequently used data.
    \item It is faster than main memory.
    \item It acts as a permanent backup of the contents of RAM.
    \item It helps exploit locality of reference.
  \end{enumerate}
\item \textit{[Original]} RAM is a \underline{\hspace{2.2cm}} memory, whereas
  ROM is \underline{\hspace{2.2cm}}.
  \begin{enumerate}[label=\Alph*.]
    \item non-volatile; volatile
    \item volatile; non-volatile
    \item secondary; primary
    \item read-write firmware; cache
  \end{enumerate}
\item \textit{[Original]} Which pairing is correctly matched?
  \begin{enumerate}[label=\Alph*.]
    \item Plotter --- input device
    \item Scanner --- output device
    \item MICR --- reads magnetic ink characters
    \item Sound card --- input device
  \end{enumerate}
\item \textit{[Original]} Evaluate the statements:
  \begin{enumerate}[label=\Roman*.]
    \item Registers are the fastest storage in a computer.
    \item Cache memory is faster than main memory.
    \item RAM is non-volatile and retains data when the power is off.
  \end{enumerate}
  \begin{enumerate}[label=\Alph*.]
    \item I only
    \item I and II only
    \item II and III only
    \item I, II, and III
  \end{enumerate}
\item \textit{[Original]} Which statement correctly distinguishes a digital
  signature from encryption?
  \begin{enumerate}[label=\Alph*.]
    \item A digital signature encrypts the body so only the recipient reads it.
    \item Encryption provides authenticity while a signature provides secrecy.
    \item A digital signature provides authenticity and integrity, while
      encryption provides confidentiality.
    \item Signatures and encryption are different names for the same process.
  \end{enumerate}
\item \textit{[Original]} Which statement correctly distinguishes POP3 from
  IMAP?
  \begin{enumerate}[label=\Alph*.]
    \item POP3 keeps mail synced on the server; IMAP deletes it locally.
    \item POP3 downloads mail to the local machine; IMAP keeps mail synced
      on the server.
    \item Both POP3 and IMAP send outgoing mail to the printer.
    \item POP3 and IMAP are different names for the same web browser.
  \end{enumerate}
\item \textit{[Original]} Assertion (A): IPv6 uses 128-bit addresses.
  Reason (R): HTTP/3 runs over QUIC rather than TCP.
  \begin{enumerate}[label=\Alph*.]
    \item Both (A) and (R) are true, and (R) is the correct explanation of (A).
    \item Both (A) and (R) are true, but (R) is not the explanation of (A).
    \item (A) is true, but (R) is false.
    \item (A) is false, but (R) is true.
  \end{enumerate}
\item \textit{[Original]} A nested IF in Excel tests a cell that is blank.
  The formula returns the value meant for the ``equal to zero'' branch.
  What does this show?
  \begin{enumerate}[label=\Alph*.]
    \item Blank input is treated exactly like a typed number in every case.
    \item Blank and non-numeric input must be traced branch by branch; the
      result depends on how the test treats the blank.
    \item Nested IF never accepts an empty cell under any condition.
    \item The equals sign at the formula start decides the branch taken.
  \end{enumerate}
\item \textit{[Original]} A quad-core 2.4 GHz machine with 16 GB RAM and a
  1 TB SSD runs one program slowly while other programs run normally. What
  should be diagnosed first?
  \begin{enumerate}[label=\Alph*.]
    \item The stated bottleneck: multithreading, memory use, or disk
      activity of that program.
    \item The CPU clock, because 2.4 GHz is always too slow.
    \item The monitor size, because display size sets program speed.
    \item The keyboard layout, because input devices set execution speed.
  \end{enumerate}
\item \textit{[Original]} Two ranges must be counted where both conditions
  hold (range 1 meets criterion 1 AND range 2 meets criterion 2). Which
  formula approach is correct?
  \begin{enumerate}[label=\Alph*.]
    \item COUNT with a single range and no criteria.
    \item COUNTIFS with range--criterion pairs for each condition.
    \item A plain SUM with no conditions attached.
    \item An AVERAGE of the two ranges with no criteria.
  \end{enumerate}
\end{enumerate}
\end{document}
